\section{Recommendations}
\label{sec:recommendations}

\begin{enumerate}
  \item \label{rec:enclosure}
    \textbf{Convert MC06 and MC07 to enclosed path analysers.} The two sites
    account for 61 percent of network wide stationarity rejections and sit
    below the network mean data return in eleven months of twelve, as
    Table~\ref{tab:monthly-record} shows. The mechanism is salt film
    accumulation on the open path windows between cleaning visits, and no
    cleaning schedule that the site access permits will fix it. An enclosed
    path instrument with a heated inlet removes the failure mode at the cost of
    a larger power budget, which both sites can carry with a 40 percent array
    expansion.
  \item \label{rec:profile}
    \textbf{Install a vertical concentration profile at one saltmarsh and one
    dune site.} Section~\ref{sec:results-advection} identifies an advective term
    of 8 to 14 percent of the midday carbon flux on sea breeze days that the
    single level eddy covariance measurement cannot resolve. A six level
    profile with a sequential sampling manifold would allow the storage and
    horizontal advection terms to be estimated directly. MC01 and MC08 are the
    natural choices because both have mains power and the longest continuous
    records.
  \item \label{rec:radiometer}
    \textbf{Shorten the radiometer calibration interval to six months.} The
    observed drift of 3.8 percent over twelve months is nearly twice the
    protocol tolerance and propagates directly into the energy balance closure
    that the network uses as its independent quality check. A six month interval
    with a travelling reference standard would hold the drift within tolerance
    at modest cost.
  \item \label{rec:threshold}
    \textbf{Report the friction velocity threshold uncertainty separately in all
    downstream products.} Table~\ref{tab:uncertainty} shows that the threshold
    contributes more uncertainty than every measurement term combined. Products
    that fold it into a single interval will mislead users into treating a
    methodological choice as a measurement limit.
  \item \label{rec:footprint}
    \textbf{Reassess the MC06 fetch or relocate the tower.} With 340 metres of
    fetch, MC06 fails the footprint criterion on 19.8 percent of records against
    a network mean of 6.2 percent. A relocation 400 metres inland along the same
    dune ridge would roughly double the fetch in the prevailing sector without
    changing the cover class.
\end{enumerate}

\section{Archive and reproducibility}
\label{sec:archive}

The archive holds the raw 20 hertz records, the half hourly fluxes at every
processing stage, and the configuration that produced them. Every archived flux
carries the identifier of the configuration used, so the effect of a
reprocessing campaign can be isolated. The configuration for the 2024 annual
release is reproduced below.

\begin{lstlisting}[caption={Processing configuration for the 2024 annual release.},label={lst:config}]
release: 2024-annual
chain_version: 4.2.1

despike:
  method: median_absolute_deviation
  threshold: 5.5
  window_minutes: 5
  max_consecutive: 4

rotation:
  method: planar_fit
  sector_width_degrees: 30
  refit_interval_days: 30

covariance:
  averaging_minutes: 30
  detrend: block_average
  lag_search_seconds: [-2.0, 2.0]

corrections:
  density: true
  spectral: empirical_transfer_function
  reference_cospectrum: site_fitted
  low_frequency: ogive

footprint:
  model: parameterised_2d
  source_fraction: 0.80
  min_target_fraction: 0.70

quality:
  stationarity_percent: 30
  integral_turbulence_percent: 30
  ustar_threshold_network: 0.13
  ustar_threshold_range: [0.08, 0.16]
  reject_precipitation: true

gapfill:
  method: marginal_distribution_sampling
  window_days: 7
  drivers: [shortwave_in, air_temperature, vapour_pressure_deficit]
  max_gap_days: 12
\end{lstlisting}

The one entry that a user is likely to want to change is
\texttt{ustar\_threshold\_network}. Reprocessing the full year at a different
threshold takes approximately four hours on a single workstation, and the
\texttt{ustar\_threshold\_range} entry records the interval over which the
sensitivity in Table~\ref{tab:uncertainty} was evaluated.

\section*{Acknowledgements}
\addcontentsline{toc}{section}{Acknowledgements}

The \network{} is funded by the coastal observation programme of the Institute
for Boundary Layer Science. Site access at MC05 and MC07 is granted by the
Roseholm Estate and at MC01 by the Kirkhavn Wildfowl Trust. Field maintenance
through 2024 was carried out by Petra Lindholm and Osmund Kaale, who kept eight
towers running through a winter that took out two of them, and the archive is
maintained by Elias Grothe.
